% Part: proof-theory
% Chapter: sequent-calculus
% Section: interpretation-rules
%READERNOTE{SOL6-A329: दोन-नियमांच्या विधानात असत्य स्थिरांकाचा अपवाद स्पष्ट केला; दुर्बलीकरणाची आवश्यकता दाखवणाऱ्या उदाहरणाला योग्य अणुरूप व नियम-संच scope दिला.}

\documentclass[../../../include/open-logic-section]{subfiles}

\begin{document}

\olfileid{pt}{seq}{irl}

\olsection{नियमांचा अर्थ}

क्रमवर्तीचा अर्थ आठवा: !!a{structure}~$\Struct{M}$ मध्ये
क्रमवर्ती~$\Gamma \Sequent \Delta$ सत्य असते, जर
किमान एक $!A \in \Gamma$ असत्य असेल किंवा किमान
एक~$!A \in \Delta$ सत्य असेल. \Log{G3c} चे
तार्किक नियम त्यांच्या प्रमुख !!{formula}s च्या
सत्यतेच्या अटी नियमबद्ध करतात असे वाचता येते.
\RightR{\lif} विचारात घ्या. त्याचे प्रमुख !!{formula}
$!A \lif !B$ आहे. $!A$ असत्य असेल किंवा~$!B$
सत्य असेल तेव्हा ते सत्य असते. म्हणून क्रमवर्ती
$\ \Sequent !A \lif !B$ ही $!A \Sequent !B$
सत्य असेल तेव्हाच सत्य असते. या क्रमवर्तींच्या
पूर्वांगात आणि उत्तरांगात संदर्भातील !!{formula}s
$\Gamma$, $\Delta$ जोडल्यास \RightR{\lif}
नियमाचे नेमके निष्कर्ष व पूर्वपक्ष मिळतात.
\LeftR{\lif} साठीही स्थिती अशीच आहे. $!A
\lif !B$ हे $!A$~सत्य आणि $!B$~असत्य असेल तेव्हा
असत्य असते. म्हणून $!A \lif !B
\Sequent \ $ ही क्रमवर्ती $\ \Sequent !A$ आणि
$!B \Sequent \ $ या दोन्ही सत्य असतील तेव्हाच
सत्य असते. \LeftR{\lif} चा निष्कर्ष त्याच्या
पूर्वपक्षांच्या संयुक्त सत्यतेशी समतुल्य आहे. या
रचनेमुळे \Log{G3c} चे नियम निर्दोष असतात (सर्व
पूर्वपक्ष सत्य असतील तर निष्कर्षही सत्य असतो).
त्यांच्यामुळे पूर्णताही मिळते.

खरे तर, कोणत्याही सत्यता-फलनीय संयोजकाच्या
सत्यतेच्या अटी अशा क्रमवर्तींच्या संचाने मांडता येतात.
कारण अभिजात तर्कशास्त्रातील प्रत्येक सत्यता-फलन
संयुक्त सामान्य रूपातील !!a{formula} ने दाखवता येते:
अण्वीय आणि निषेधित अण्वीय !!{formula}s यांच्या
वियोजनांचे संयोजन. उदाहरणार्थ, $!A
\oplus !B$ हे $!A$ आणि $!B$ यांपैकी नेमके एकच
सत्य असेल तेव्हाच सत्य मानू; यालाच ``अनन्य किंवा''
म्हणतात. मग $!A \oplus !B$ हे $(!A \lor !B) \land
(\lnot !A \lor \lnot !B)$ सत्य असेल तेव्हाच सत्य
असते; आणि $(\lnot !A \lor !B)
\land (!A \lor \lnot !B)$ सत्य असेल तेव्हा ते असत्य
असते. म्हणून $\oplus$ साठी क्रमवर्ती-कलनाचे नियम
पुढीलप्रमाणे देता येतील:
\begin{align*}
&
\Axiom$\Gamma \fCenter \Delta, !A, !B$
\Axiom$!A, !B, \Gamma \fCenter \Delta$
\RightLabel{\RightR{\oplus}}
\BinaryInf$\Gamma \fCenter \Delta, !A \oplus !B$
\DisplayProof
\\
& \Axiom$!A, \Gamma \fCenter \Delta, !B$
\Axiom$!B, \Gamma \fCenter \Delta, !A$
\RightLabel{\LeftR{\oplus}}
\BinaryInf$!A \oplus !B, \Gamma \fCenter \Delta$
\DisplayProof
\end{align*}
हे नियमही निर्दोष व पूर्ण असतील. त्यांचे निष्कर्ष सर्व
पूर्वपक्ष सत्य असतील तेव्हाच सत्य असतात.

\begin{prob}
$!A \downarrow !B$ हे $!A$ किंवा~$!B$ यांपैकी
एकही सत्य नसल्यास आणि तरच सत्य आहे असे माना
(``NOR''). त्याचे \Log{G3c} मधील नियम कोणते असतील?
($!A \downarrow !B$ सत्य असेल तेव्हाच एकाच वेळी
सत्य असणाऱ्या दोन क्रमवर्ती शोधा. त्यातून
\RightR{\downarrow} नियम मिळेल. त्यानंतर $!A
\downarrow !B$ असत्य असेल तेव्हाच सत्य असणारी एक
क्रमवर्ती शोधा. त्यातून \LeftR{\downarrow} नियम मिळेल.)
\end{prob}

\Log{G3c} मध्ये प्रत्येक संयोजक आणि संख्यक
यांचे नेमके दोन नियम आहेत: एक डावीकडचा आणि एक
उजवीकडचा. येथे~$\lfalse$ हे स्वतंत्र तार्किक स्थिरांक
आहे; त्यासाठी डावीकडील स्वयंसिद्धक असते, उजवीकडील
स्वतंत्र नियम नसतो. पण \Log{G1c} मध्ये दोन
\LeftR{\land} आणि दोन \RightR{\lor} नियम आहेत.
याचे कारण असे की \Log{G1c} चे प्रतिमान असलेल्या
गेन्ट्झेनच्या मूळ~\Log{LK} पद्धतीला एक महत्त्वाच्या
अभिजातेतर तर्कशास्त्रासाठी, म्हणजे अंतःप्रज्ञावादी
तर्कशास्त्रासाठी,~\Log{LJ} हा प्रकार होता.
अंतःप्रज्ञावादी तर्कशास्त्रासाठी निर्दोष व पूर्ण
पद्धत मिळवण्याकरिता~\Log{LJ} मधील प्रत्येक
क्रमवर्तीच्या उत्तरांगात जास्तीत जास्त एकच
!!{formula} ठेवण्याचे बंधन आहे. त्यामुळे
\Log{G3c} चा \RightR{\lor} नियम वापरता येत
नाही, कारण त्याच्या पूर्वपक्षाच्या उत्तरांगात $!A$
आणि~$!B$ ही दोन्ही येतात. म्हणून गेन्ट्झेनने
\Log{LJ} साठी \RightR{\lor} नियमांची जोडी
वापरली, आणि तीच~\Log{LK} साठीही वापरली.
तदनुरूप, अंतःप्रज्ञावादी प्रकार~\Log{G1i} हा
प्रत्येक उत्तरांगात जास्तीत जास्त एकच !!{formula}
ठेवण्याचे बंधन घातलेली~\Log{G1c} पद्धत आहे.
(\Log{G3c} चाही अंतःप्रज्ञावादी प्रकार आहे;
मात्र त्यातील काही नियम \Log{G3c} मधील संबंधित
नियमांहून वेगळे आहेत. उदाहरणार्थ, त्यातही
\RightR{\lor} नियमांची जोडी वापरतात.)

\Log{G1c} चे संरचनात्मक नियम निर्दोष आहेत
हे सहज दिसते. उदाहरणार्थ, $\Gamma \Sequent \Delta$
च्या डावीकडे असत्य !!{formula} किंवा उजवीकडे सत्य
!!{formula} असेल, तर $\Gamma \Sequent
\Delta, !A$ मध्येही तेच असते; $!A$ सत्य आहे की
असत्य, याने फरक पडत नाही. म्हणून
\RightR{\Weakening} चा पूर्वपक्ष सत्य असेल तर
निष्कर्षही सत्य असतो. येथे उलट दिशेचा दावा मात्र
सत्य नाही: निष्कर्ष $!A$ सत्य असल्यामुळे सत्य असू
शकतो; तेव्हा पूर्वपक्ष सत्य असेल याची खात्री नाही.

संरचनात्मक नियम मुळात का हवेत? उदाहरणार्थ,
$\ \Sequent !A \lor \lnot !A$ !!{prove} करण्यासाठी
आकुंचन (\RightR{\Contraction},
\LeftR{\Contraction}) लागते. $!A \Sequent !A$
पासून सुरुवात केल्यास, आधी~\RightR{\lnot} वापरून
$\Sequent !A, \lnot !A$ मिळते. मग
\RightR{\lor} च्या दोन आवृत्त्या क्रमाने वापरता येतात:
एकदा~$\lnot !A$ पासून $!A \lor \lnot
!A$ मिळवण्यासाठी आणि एकदा~$!A$ पासून
$!A \lor \lnot !A$ मिळवण्यासाठी. त्यातून $\ \Sequent !A \lor \lnot
!A, !A \lor \lnot !A$ मिळते. म्हणून \Log{G1c}
मध्ये $\Sequent !A \lor \lnot !A$ मिळवण्यासाठी
!!a{proof} मधील त्याच !!{formula} च्या दोन
उपस्थितींचे ``आकुंचन'' करता आले पाहिजे:
\[
\Axiom$!A \fCenter !A$
\RightLabel{\RightR{\lnot}}
\UnaryInf$\fCenter !A, \lnot !A$
\RightLabel{\RightR{\lor}}
\UnaryInf$\fCenter !A, !A \lor \lnot !A$
\RightLabel{\RightR{\lor}}
\UnaryInf$\fCenter !A \lor \lnot !A, !A \lor \lnot !A$
\RightLabel{\RightR{\Contraction}}
\UnaryInf$\fCenter !A \lor \lnot !A$
\DisplayProof
\]
खरे तर, दुर्बलीकरण किंवा आकुंचन नसताना \Log{G1c}
मध्ये !!{prove} न करता येणाऱ्या वैध क्रमवर्ती आहेत.
उदाहरणार्थ,~$!A$ आणि~$!B$ ही वेगळी अणुरूप
सूत्रे घ्या. दुर्बलीकरण व आकुंचन वगळलेल्या तक्त्यातील
\Log{G1c} नियमांनी क्रमवर्ती
$!A \Sequent !B \lif !A$ ची !!a{proof}
\RightR{\lif} नेच संपली पाहिजे (कारण फक्त
$!B \lif !A$ प्रमुख असू शकते); त्यासाठी
$!B, !A \Sequent !A$ हा पूर्वपक्ष हवा. तो
$!A \Sequent !A$ या स्वयंसिद्धकापासून मिळवायला
\LeftR{\Weakening} लागते:
\[
\Axiom$!A \fCenter !A$
\RightLabel{\LeftR{\Weakening}}
\UnaryInf$!B, !A \fCenter !A$
\RightLabel{\RightR{\lif}}
\UnaryInf$!A \fCenter !B \lif !A$
\DisplayProof
\]

\Log{G3c} मध्ये आकुंचन आणि दुर्बलीकरण
हे नियम अजिबात लागत नाहीत. दुर्बलीकरण आधीच
स्वयंसिद्धकांत अंतर्भूत आहे. दोन पूर्वपक्ष असलेल्या
नियमांमधील संदर्भ~$\Gamma$, $\Delta$ यांच्या दोन
प्रती एकत्र करून आकुंचन बहुतेक ठिकाणी टाळता येते.
\Log{G1c} आणि \Log{G3c} या दोन्हींमध्ये असे
``संदर्भ सामायिक करणारे'' नियम आहेत. एकच पूर्वपक्ष
असलेल्या नियमांत आकुंचन फक्त \RightR{\lor},
\LeftR{\land}, \RightR{\lexists} आणि~\LeftR{\lforall}
यांसाठी लागते. \Log{G3c} मधील या नियमांच्या
आवृत्त्या आकुंचन पूर्णपणे अनावश्यक करतात.
\Log{G2c} ही आणखी एक पद्धत आहे (
\olref{tab:G2c} पाहा); तिच्या दोन-पूर्वपक्षी नियमांना
स्वतंत्र संदर्भ असतात. \Log{G2c} मध्ये आकुंचन
\Log{G1c} पेक्षाही अधिक महत्त्वाचे आहे.

\subfile{rules-G2c}

\end{document}
